\section{Related Work}
\label{sec:related_work}
\label{sec:related_sim}

AEGIS-Q separates storage protection, GPU/storage placement, and PQC freshness; it is not a stronger deployed encryption layer.  fscrypt and dm-crypt are mature kernel baselines~\cite{fscrypt, dmcrypt}; gocryptfs is the closest user-space encrypted-filesystem comparison point~\cite{gocryptfs}.  fs-verity and dm-integrity validate different integrity boundaries, so this revision treats plaintext/gocryptfs warm-cache rows as current evidence and leaves a full fscrypt/dm-crypt matrix open.  AEGIS-Q's distinct artifact is the combination of a per-file authenticated FUSE format, storage-visible latency/elastic policy, optional GPU maintenance with CPU fallback, and externally anchored replay checks in one mounted prototype.

Prior GPU/storage systems validate different assumptions.  FlashNeuron and Fastensor use SSD staging for learning workloads~\cite{flashneuron, fastensor}; SnuQS and InfScaler motivate managed storage capacity~\cite{snuqs_storage, infscaler}; SnuQS alone is not a validation of this filesystem format~\cite{snuqs_storage}.  GPUfs, SPDK, and GPUDirect Storage need API and hardware evidence beyond ordinary FUSE \texttt{pread}/\texttt{pwrite}.  Edge pressure is real~\cite{breaking_edge}, and FlashNeuron illustrates why storage-backed execution needs separate foreground evidence~\cite{breaking_edge, flashneuron}.  Batch GPU work appears under datacenter assumptions~\cite{flashneuron, fscrypt_gpu, fastensor}, including Fastensor/fscrypt-GPU comparisons~\cite{fastensor, fscrypt_gpu}; fscrypt-GPU itself is not a matched baseline for this prototype~\cite{fscrypt_gpu}.  OP-TEE and dm-crypt use different protection boundaries~\cite{optee, dmcrypt}.  Thus the portable idea is measured admission, while CUDA/cuPQC evidence is platform-specific.

ML-KEM establishes shared secrets rather than per-block bulk encryption~\cite{mlkem}, so AEGIS-Q keeps AES-GCM in the data plane and uses ML-KEM only for the optional batched key-plane experiment.  OP-TEE provides a different trust boundary~\cite{optee}; it does not make a replayable file witness fresh.  A TPM freshness root still needs explicit NV authorization, key-release policy, and recovery semantics beyond the present measurements.
